\chapter{Verification and Validation}
\label{cap:validation}

This chapter reports how the platform was checked for correctness
(\textit{verification} --- ``are we building the product right?'') and how
it was confirmed to address the problems identified during discovery
(\textit{validation} --- ``are we building the right product?''). The
verification effort rests on three complementary layers: static analysis
through the TypeScript compiler, an automated unit-test suite executed on
every change, and manual acceptance testing with the client. Validation is
addressed through a structured survey mapped one-to-one onto the pain points
of Chapter~\ref{cap:introduction}, and the chapter closes with the security
review that accompanied the project.

%----------------------------------------------------------------------------------------

\section{Verification Strategy}
\label{sec:val-verification}

Verification follows the same boundaries as the \gls{onion} described in
Chapter~\ref{cap:design}. Each layer is checked with the technique best
suited to it: the type system guarantees structural correctness across all
layers; pure Domain and Application logic is exercised by fast, isolated
unit tests; and the Presentation boundary (authentication and route gating)
is tested through request-level assertions. Because every external
dependency is reached through a port, the suite mocks at those interfaces
and never touches the network or the database, which keeps the whole run
deterministic and fast enough to execute on every push.

\subsection{Static Analysis and Type Safety}
\label{sec:val-static}

The first line of defence is the TypeScript compiler running in
\texttt{strict} mode. The command \texttt{tsc -{}-noEmit} type-checks the
whole codebase without producing output and is the canonical local
validation step for this project, in preference to a full production build
(the build intentionally fails locally because it depends on
deployment-only environment variables for the interview routes). Strict
mode forbids implicit \texttt{any}, enforces null-safety, and makes the
path aliases (\texttt{@server}, \texttt{@client}, \texttt{@/lib}) resolve
identically in source and in tests.

A code audit conducted in May~2026 found that the weakest point of the type
guarantees was the database boundary, where rows were cast through
\texttt{camelizeKeys<any>} and lost their types exactly where data crosses
from PostgreSQL into the Domain. This observation is recorded as a
maintainability item; it illustrates that type checking is necessary but
not sufficient, and motivates the unit tests that pin the behaviour of the
conversion utilities and the schema validators.

\subsection{Unit Tests}
\label{sec:val-unit}

\subsubsection{Tooling}

The suite runs on \textbf{Vitest~4.1.5} with the V8 coverage provider. Vitest
was chosen over Jest because it natively supports ECMAScript modules (no
CommonJS transform step), reuses the same path aliases as the application
through its Vite resolver, and starts quickly thanks to the esbuild
transformer. Tests live in the top-level \texttt{tests/} directory and
alongside the code under \texttt{src/}, and are executed with
\texttt{vitest run} (single pass, used in CI) or \texttt{vitest -{}-watch}
during development. Coverage is produced with \texttt{vitest run
-{}-coverage} and uploaded as a continuous-integration artefact.

\subsubsection{Suite Composition}

At the time of writing the suite comprises \textbf{18 test files and 305
test cases, all passing}. Table~\ref{tab:test-inventory} lists every file
with its case count, the application layer it targets, and the concern it
covers. The suite deliberately concentrates on the layers where logic
lives --- the Domain (pure decision functions) and the Application
(orchestration) --- while the thin Infrastructure adapters are verified
indirectly through the use cases that depend on them.

\begin{table}[h]
\centering
\caption[Automated test inventory]{Automated test inventory at submission
(18 files, 305 cases, all passing).}
\label{tab:test-inventory}
\small
\begin{tabular}{p{0.30\linewidth}r p{0.18\linewidth}p{0.32\linewidth}}
\toprule
\textbf{File} & \textbf{Cases} & \textbf{Layer} & \textbf{Concern} \\
\midrule
\texttt{job-fit} & 47 & Domain & Job-anchored matching engine \\
\texttt{middleware-auth} & 38 & Presentation & Route gating (401/403) \\
\texttt{interview-runtime} & 37 & App + Infra & Interview rubric + evidence \\
\texttt{job-matching-bridge} & 25 & Application & DB-row $\rightarrow$ matcher bridge \\
\texttt{cv-validation} & 19 & Application & CV extraction schema \\
\texttt{upload-use-cases} & 18 & Application & CV upload pipeline \\
\texttt{analytics-catalog} & 16 & Domain & Widget-spec + injection guard \\
\texttt{notifications-route-auth} & 16 & Presentation & API route authorisation \\
\texttt{scoring} & 13 & Domain & Scoring + CEFR estimation \\
\texttt{adidas-job-scraper} & 10 & Application & Job-listing scraper \\
\texttt{logger-redaction} & 10 & Infrastructure & PII redaction in logs \\
\texttt{text-extraction} & 10 & Infrastructure & PDF/DOCX extraction \\
\texttt{escape-or-term} & 9 & Infrastructure & PostgREST \texttt{.or()} escaping \\
\texttt{job-requirements-schema} & 9 & Application & JD extractor schema \\
\texttt{listing-posted-date} & 9 & Domain & Fuzzy posted-date parsing \\
\texttt{job-shortlist-use-cases} & 8 & Application & Per-job shortlist \\
\texttt{interview-token} & 6 & Infrastructure & HMAC token round-trip \\
\texttt{cv-fields-of-work} & 5 & Application & Field-of-work tagging \\
\midrule
\textbf{Total} & \textbf{305} & & \\
\bottomrule
\end{tabular}
\end{table}

\subsubsection{What is Tested, and Why}

The suite is not uniform by design; the depth of testing in each area is
proportional to its risk and to how much consequential logic it contains.

\paragraph{Decision-critical Domain logic (highest priority).}
The job-anchored matching engine (\texttt{computeJobFit}) is the most
heavily tested unit, with 47 cases. Because hiring decisions must be
explainable rather than opaque, the matcher is a pure function with no
external dependencies, and its tests cover every one of the seven criteria
(field, experience-in-field, seniority, required and preferred skills,
languages, education), the ``average of applicable criteria only'' rule,
and the \texttt{isEligible} flag computed as the conjunction of applicable
\texttt{met} flags. A dedicated \textit{golden skill-match corpus} locks in
expected match and no-match outcomes for realistic skill pairs --- for
example that \texttt{js} matches \texttt{javascript} and \texttt{k8s}
matches \texttt{kubernetes} through curated synonym groups and phrase
aliases, while genuinely distinct products such as AWS and Azure, or Java
and JavaScript, are kept apart. This corpus exists specifically to prevent
silent regressions in matching accuracy when the synonym tables are edited.
The scoring service (13 cases) is tested in the same spirit, covering score
ranges, relative ordering, boundary capping, and the keyword-to-\gls{cefr}
mapping.

\paragraph{The glue between layers.}
The \texttt{job-matching-bridge} suite (25 cases) was added after a
production defect in which the bridge that converts a database row into the
matcher's input read the wrong property (\texttt{title} instead of
\texttt{jobTitle}), silently emptying the evidence text that supports
title-based seniority matching. The pure matcher was already well tested,
yet the untested ``glue'' in the Application layer carried the bug. The
new suite tests the loose-date parser, the experience-duration calculation,
the requirement-building helper, and includes an explicit regression case
asserting that the job title is carried through. The lesson --- that pure
logic and the code wiring it together must both be tested --- is revisited
in Chapter~\ref{cap:conclusions}.

\paragraph{Data-integrity boundaries.}
The platform ingests data from Large Language Models, whose output cannot be
trusted structurally. Four schema suites (\texttt{cv-validation},
\texttt{job-requirements-schema}, \texttt{cv-fields-of-work}, and the
spec validation inside \texttt{analytics-catalog}) assert that
Zod schemas accept well-formed payloads, reject missing required fields,
normalise values such as URLs and e-mail addresses, and tolerantly discard
model-invented values that fall outside the canonical vocabulary (for
instance, fields of work outside the sixteen allowed).

\paragraph{Security-sensitive paths.}
Several suites exist purely to defend security and compliance properties:
\texttt{middleware-auth} (38 cases) and \texttt{notifications-route-auth}
(16 cases) assert that unauthenticated callers receive \texttt{401} and
non-HR callers receive \texttt{403} on protected routes; \texttt{escape-or-term}
(9 cases) verifies that user search terms are escaped before reaching
PostgREST \texttt{.or()} filters, closing an injection vector found in the
audit; \texttt{logger-redaction} (10 cases) confirms that e-mail addresses
and tokens are masked in structured logs; and the analytics catalog tests
confirm that the strict widget-specification schema rejects injected keys,
so that the constrained chart builder can never be coerced into running an
arbitrary query.

\paragraph{Module ownership note.}
Two suites, \texttt{interview-runtime} (37 cases) and \texttt{interview-token}
(6 cases), exercise the AI Skill Interviewer integration. The conversational
FastAPI service itself was developed by the second team member; the tests
listed here cover the Next.js side of that integration --- session creation,
HMAC token signing and expiry, transcript-turn persistence, and the
rubric-validation guardrails applied before an evaluation is stored. They
are part of the shared suite but are attributed accordingly to keep the
contribution boundary explicit.

\subsubsection{Test Architecture and Mocking}

The tests mirror the application's layering. Because every external
dependency is reached through a Domain port, each port has a hand-written
mock factory (for example a fake candidate repository, CV parser, storage
service, and text extractor for the upload pipeline). External libraries
that have no port --- the PDF and DOCX parsers, and the Supabase client used
by the interview routes --- are replaced with module mocks. No mocking
library is used; plain stub functions with explicit return values keep each
test file self-contained, with its own fixtures and no shared global setup.
This makes failures local and easy to read, at the small cost of some
duplication between files.

\subsection{Manual Acceptance Testing with the Client}
\label{sec:val-manual}

Automated tests cannot judge whether a screen \textit{feels} faster to a
recruiter, so the automated suite was complemented by structured acceptance
sessions with the client. Working from a seeded dataset of representative
CVs and job descriptions, the team walked through the end-to-end recruiter
journey --- uploading and parsing a CV, opening a job and ranking candidates
against it, shortlisting, contacting a candidate, and reviewing analytics ---
and recorded the client's feedback against the original pain points. These
sessions were also the principal vehicle for sprint reviews, doubling as
both acceptance testing and Scrum ceremony (Section~\ref{sec:analysis-pm}).

%----------------------------------------------------------------------------------------

\section{Validation through Client Survey}
\label{sec:val-survey}

Verification establishes that the software behaves as specified; validation
asks whether the specification actually solves the client's problem. To
answer this, the nine pain points of Table~\ref{tab:pain-points-overview}
were turned into a survey instrument with one question per pain point, so
that every reported friction could be traced to an explicit measure of
whether the platform relieved it. The survey captured both the perceived
severity of each pain point before the tool and the expected improvement
with it, distributed to the HR stakeholders involved in discovery.

The analysis focuses on the three highest-priority pain points identified
during discovery: the time spent extracting information from CVs manually
(pain point~5), the absence of first-pass language and skill screening
before an interview (pain point~3), and the lack of a structured process to
keep promising-but-rejected candidates engaged (pain point~9). The
corresponding response distributions are reproduced as bar charts in this
section. The validity of the findings is bounded by a small sample drawn
from a single organisation; the results are therefore read as directional
confirmation that the prioritisation was correct, not as a statistically
generalisable claim. The full instrument and the raw responses are
reproduced in Appendix~\ref{app:survey}.

%----------------------------------------------------------------------------------------

\section{Security Review}
\label{sec:val-security}

Security was treated as a first-class verification concern rather than an
afterthought, for two reasons: the platform stores special-category
personal data under \gls{gdpr}, and it disables PostgreSQL \gls{rls} in
favour of enforcing every access decision in server-side code reached
through the service role. That design choice concentrates the
authorisation burden in the application, which is precisely why the
authentication and route-gating suites described above are central to the
security posture.

A deep code audit in May~2026 produced a prioritised findings log that
drove concrete remediation. The most serious item was a leaked service-role
credential committed in an obsolete migration script; the file was removed
and the key rotated in the Supabase dashboard, with the corresponding
environment variables updated across all deployment targets. The audit also
identified --- and the team resolved --- a PostgREST \texttt{.or()} injection
vector in the search filters, fixed by introducing and unit-testing the
\texttt{escapeOrTerm} helper and applying it at every call site. Remaining
items, such as extracting the interview routes behind a use-case layer and
replacing residual \texttt{any} casts at the repository boundary with typed
row interfaces, are recorded as tracked maintainability work rather than
open vulnerabilities.

Finally, a sweep with the Supabase Security Advisor was used to confirm the
intended state of the database surface. The remaining advisor entries
correspond to deliberate design decisions --- chiefly that \gls{rls} is
disabled by design on internal tables that are never reached by an
end-user session --- and each is recorded with its rationale in a decisions
log so that an intentional choice is never mistaken for an oversight.

%----------------------------------------------------------------------------------------

\section{Performance Considerations}
\label{sec:val-performance}

Performance was addressed pragmatically, at the points the data model makes
hot. Indexes back the read paths that the recruiter workflow exercises most
often: the \texttt{job\_matches} and \texttt{job\_shortlists} tables for the
per-job ranking and shortlist screens, and a GIN index on
\texttt{experiences.fields\_of\_work} so that candidates can be filtered by
field of work without a full scan. Expensive computation is cached rather
than repeated: a job description is parsed by the Large Language Model only
on the first time an HR user opens the ranking screen for that job, and the
result is persisted with a schema version so that it is reused until the
source changes; likewise the top-ranked candidates for a job are persisted
as a \texttt{job\_matches} cache rather than recomputed on every view.

The principal known bottleneck is the candidate fan-out inside
\texttt{matchCandidatesToJob}, which currently loads the candidate pool and
evaluates the matcher in application code. The remediation identified during
the audit --- pre-filtering candidates in SQL through the fields-of-work GIN
index so that only plausibly relevant profiles reach the matcher --- is
recorded as planned work; the existing matcher tests are written so that
they continue to hold once that optimisation is introduced.

%----------------------------------------------------------------------------------------

\section{Limitations Found During Validation}
\label{sec:val-limitations}

Validation also surfaced limitations that are reported here honestly rather
than hidden. The browser-based speech-to-text and text-to-speech used by the
assessment experience rely on Web Speech APIs available only in
Chromium-based browsers (Chrome and Edge), so the spoken-interaction path
degrades on other browsers. The Large Language Model usage runs against a
provider free tier, which imposes a practical cost-and-rate ceiling on bulk
CV parsing and is mitigated, but not removed, by the Groq-to-OpenAI fallback
and the lazy parsing of job descriptions. Some \gls{gdpr} obligations are
still operationally manual: the deletion-aware data model is in place, but
the right-to-erasure routine is executed as a script rather than as a
self-service flow. Finally, the constrained analytics builder, while safe by
construction, exposes a deliberately narrow catalogue of metrics and
dimensions, so some questions an analyst might ask cannot yet be expressed.
These items frame the future work discussed in
Chapter~\ref{cap:conclusions}.
